% ============================================================================
% Chapitre 10 — Conclusion générale
% ============================================================================
\chapter{Conclusion générale}

\section{Synthèse du travail réalisé}

Ce projet a donné naissance à \textbf{\caimp{}}, une plateforme de supervision
d'infrastructure open source qui repose sur une philosophie inédite dans le
monde des outils de monitoring : répondre directement à la question
« \textit{Pourquoi mon serveur est-il cassé ?} » plutôt que de simplement
afficher des graphes que l'opérateur devra interpréter seul.

En l'espace d'un mois, une équipe de trois étudiants ingénieurs a conçu,
développé, testé et documenté une plateforme complète de 19 microservices,
plus de 52~000 lignes de code réparties en Python 3.12, Go 1.22 et TypeScript,
une base de données relationnelle de 22 tables intégrant TimescaleDB et pgvector,
et une interface utilisateur moderne orientée réponses. L'intégration d'un grand
modèle de langage local (Llama 3.1 8B via Ollama), d'un moteur de corrélation
de changements, et d'algorithmes de prévision statistique (Holt-Winters) place
\caimp{} à l'intersection de l'ingénierie logicielle, de la science des données
et de l'intelligence artificielle appliquée.

Les fonctionnalités clés livrées incluent : la détection d'anomalies en temps
réel par trois algorithmes complémentaires, l'explication en langage naturel
par \textsc{llm} local, la corrélation automatique des changements système avec
les incidents, la prévision sur 24 heures avec intervalles de confiance, la
déduplication intelligente des alertes, la génération de rapports PDF accessibles
aux non-techniciens, et une architecture microservices sécurisée par mTLS, JWT
et RLS PostgreSQL.

\section{Apports académiques et professionnels}

Ce projet a mobilisé un ensemble de compétences transversales rarement toutes
sollicitées dans un seul module académique :

\begin{itemize}
  \item \textbf{Architecture distribuée} : conception d'un système en microservices
    communicants via un bus d'événements NATS JetStream, avec gestion du cycle
    de vie des messages, de la durabilité et de la consommation distribuée ;
  \item \textbf{Bases de données avancées} : utilisation experte de PostgreSQL
    avec TimescaleDB (hypertables, compression columnar, politiques de rétention)
    et pgvector (index HNSW, similarité cosinus) ;
  \item \textbf{Intelligence artificielle appliquée} : intégration d'un LLM local,
    ingénierie de prompts pour un public spécifique, pipeline RAG complet ;
  \item \textbf{Ingénierie des données} : modélisation statistique (Holt-Winters),
    traitement de séries temporelles, scoring de corrélation ;
  \item \textbf{Sécurité} : conception et implémentation d'une infrastructure
    PKI interne, mTLS, JWT, RBAC, RLS --- couverture de l'ensemble de la pile
    de sécurité ;
  \item \textbf{DevOps} : containerisation complète, CI/CD, migrations versionées,
    monitoring de la plateforme elle-même.
\end{itemize}

\section{Bilan personnel}

\paragraph{CHAHIDY Rim}
Ce projet m'a permis de consolider mes compétences en développement frontend
React et de découvrir la complexité fascinante de la conception de systèmes
distribués. La réalisation du Telemetry Writer en Go a été un défi technique
stimulant qui m'a initiée aux mécanismes de concurrence natifs du langage.
J'emporte de cette expérience la conviction que la qualité d'une interface
utilisateur peut transformer un outil fonctionnel en une solution réellement
utile.

\paragraph{SAMDI Wiam}
Le travail sur le pipeline RAG et les prompts LLM a été la partie la plus
intellectuellement enrichissante de ce projet. Apprendre à « parler » à un
modèle de langage de manière à obtenir des réponses structurées et fiables est
un art qui demande de la rigueur et beaucoup de tests. Ce projet m'a aussi
confirmé l'importance des tests d'intégration dans les systèmes distribués :
les bugs les plus difficiles à trouver n'apparaissent pas dans les tests
unitaires.

\paragraph{TAÏK Youssef}
La conception de la sécurité de bout en bout --- de l'autorité de certification
interne jusqu'aux politiques RLS PostgreSQL --- a été l'aspect le plus
formateur de ma contribution. Ce projet m'a appris qu'une bonne architecture
de sécurité n'est pas une couche qu'on ajoute à la fin, mais une colonne
vertébrale qu'on construit dès le premier jour. La vision de \caimp{} comme
ingénieur SRE virtuel pour les petites équipes me semble être une direction
pleine d'avenir.

\section{Mot de la fin}

\caimp{} est né d'une conviction simple : les développeurs qui gèrent leurs
propres serveurs ne devraient pas avoir à devenir des experts en monitoring
pour savoir si leurs systèmes fonctionnent. La technologie est aujourd'hui
suffisamment mature --- LLMs locaux, bases de données vectorielles, architectures
événementielles légères --- pour automatiser l'essentiel du travail d'investigation
qu'une alerte de supervision requiert habituellement.

Nous espérons que ce travail contribuera, modestement, à démocratiser l'accès
à une supervision intelligente, et qu'il constituera une base solide pour des
équipes qui souhaiteront l'étendre et l'améliorer.

\bigskip
\begin{flushright}
\textit{Settat, 2026}
\end{flushright}
